\section{Related work and evidence scope}
\label{sec:related}
\subsection{Where numerical consistency is enforced}
Query-axis parallelism is established attention structure. FlashAttention uses tiled attention to reduce memory traffic~\cite{flashattention}; FlashAttention-2 distributes query rows across thread blocks and warps while sharing K/V~\cite{flashattention2}. These methods optimize fused evaluation and work partitioning. QB instead expresses independent queries through views of existing non-flash operators and integrates their one-query organization with phase selection, graph reuse and speculative state. Its qualification target is the output of a specified deployed numerical stack.

Table~\ref{tab:relatedmechanisms} compares the closest control points. Batch-invariant inference stabilizes reductions in normalization, matrix multiplication and attention~\cite{thinkingmachines}. CoRun combines isolated prefill, fixed-shape decode and selective kernel constraints~\cite{corun}. LLM-42 runs a fast path, verifies fixed-shape windows and repairs the corresponding KV entries~\cite{llm42}. These systems define reproducible execution contracts; matching an already deployed quantized reference is a separate requirement.

\begin{table*}[t]\centering\small
\caption{Numerical-consistency control points. Performance comparisons in this paper use GS and PQ in the same runtime.}\label{tab:relatedmechanisms}
\begin{tabularx}{\textwidth}{@{}>{\raggedright\arraybackslash}p{2.3cm}>{\raggedright\arraybackslash}p{3.3cm}YY@{}}\toprule
Approach & Consistency target & Constraint and arithmetic & State and verification\\\midrule
Batch-invariant kernels~\cite{thinkingmachines}
& Each token's result is invariant to batching and slicing
& Stable reduction/tiling rules in kernels, including attention split size
& Consistent KV layout before attention; no added verify--rollback loop\\
CoRun~\cite{corun}
& Stable token sequence across co-runners, batch positions and arrival orders in a fixed stack
& Isolated prefill and fixed-shape CUDA-Graph decode; disable decode Split-KV and constrain collectives
& Request-local cache and RNG state; shape control supplies the execution contract\\
LLM-42~\cite{llm42}
& Tokens agree with a fixed-shape verification schedule
& Dynamic fast path plus fixed-shape replay; mostly existing kernels
& Verify--rollback and replacement of decoded KV entries by verified entries\\
Phase-scoped QB
& Complete output agrees with the specified ordinary reference on the qualified panel
& Q/mask views preserve local one-query attention; surrounding computation remains batched
& Phase/layout-aware reuse, producer-state qualification and native speculative reconciliation\\\bottomrule
\end{tabularx}
\end{table*}

QB's method is the integrated graph/runtime transformation. It reuses attention operators while changing query/mask representation, preserving head semantics and binding the layout to phases and graph identity. The Vulkan AV-preserving port shows that backend admission can still require an implementation change. Local numerical correspondence and stateful endpoint qualification test different parts of this contract.

Other approaches act at different boundaries. LayerCast separates storage and computation precision~\cite{layercast}; TBIK aligns reductions across tensor-parallel configurations~\cite{tbik}; MarginGate selects verification using decision margins~\cite{margingate}. Vosti specifies engine/kernel behavior across batching, chunking and cache changes and proves guarantees under its assumptions~\cite{vosti}. Our logical mapping proof, runtime guards, operator replay and endpoint tests supply finite qualification against the specified reference.

\subsection{State and request cost}
Stage-Replay demonstrates suffix sensitivity to KV transplantation~\cite{stagereplay}, and cache-dependent divergence is reported in quantized serving~\cite{cachestate}. Our fixed-precision interventions add component-substitution counterexamples and common-prefix score comparisons. Activation-patching studies establish the importance of intervention metrics and nonadditive component effects~\cite{patching,mediators}; our four-cell score contrast tests nonadditivity beyond an endpoint pattern caused only by thresholding.

Ragged-batch synchronization errors supply a separate source of speculative failures~\cite{forensics}. We check semantic indexing, finite-precision correspondence and backend admission separately. Work on verification cost and actual execution behavior~\cite{illusion,lossless} also motivates complete-request measurement. Our matched comparisons add the repair--draft interaction and isolate QB's incremental cost relative to PQ.

\subsection{Scope of the empirical results}
The CUDA request results cover one sequence, non-flash attention, $q\leq4$, Q8\_0 weights, F16 KV and context1024. E2 estimates variation across six sessions on one host for a fixed 12-input panel; E1 samples one prompt per family at physical key extent 256 and short visible contexts. The new-input study extends formatted prompts to 811 tokens and observed continuation visibility to 938 keys. Its short mathematical and longer programming inputs vary domain and length together~\cite{rtxtransferrecords}. The Llama counterexample in Section~\ref{sec:transfer} bounds reference transfer. All three families reject the tested native-fast comparator, so qualified cost comparisons remain against the declared serial and repaired paths.

The Vulkan state interventions concern four exposed inputs and finite greedy continuations. The stopped Llama Q6\_K pilot remains unqualified. The Qwen3 Q8\_0 port has its own reference and runtime; native replay covers four complete requests, a real layer-0/key512/$q=2$ fixture and its first-query $q=1$ slice~\cite{b570nativereproduction}. The historical $k=1$ service and prefill-submission screens, including the latter's incomplete helper-process accounting, are retained in \sref{supp:service}.

The matched Qwen3 study estimates a joint draft-length/lifecycle contrast on twelve exposed inputs. Its six process sessions occur on one date; the four cluster-bootstrap intervals have nominal family coverage and quantify process-session variation on that host~\cite{rtxnativebridge}. Individual lifecycle contrasts include parity. The service estimates include the full declared request window, while CUDA Graph capture/replay mechanisms and the remaining differences between historical drivers are unidentified. Contexts beyond 1024, concurrent throughput, energy and hardware efficiency are unmeasured. Native rebuild/replay establishes source-host execution; external access to bulk tensors and pinned numerical dependencies is specified separately from saved-data reconstruction and CPU checks in \sref{supp:reproduction}.
